\documentclass[11pt]{article}
\usepackage[utf8]{inputenc}
\usepackage[T1]{fontenc}
\usepackage{lmodern}
\usepackage[spanish]{babel}
\usepackage{amsmath,amssymb}
\usepackage{graphicx}
\usepackage[left=2cm,right=2cm,top=2cm,bottom=2cm]{geometry}
\usepackage{enumitem}
\usepackage{booktabs}
\usepackage{array}
\usepackage{xcolor}
\usepackage{framed}

\spanishdecimal{.}

\setlist[enumerate,1]{label=\textbf{\arabic*.}, leftmargin=*, itemsep=6pt}
\setlist[enumerate,2]{label=(\alph*), itemsep=4pt}
\setlength{\parindent}{0pt}
\setlength{\parskip}{4pt}
\pagestyle{plain}

\newcommand{\Ein}{E_{\text{in}}}
\newcommand{\Eout}{E_{\text{out}}}
\newcommand{\Hyp}{\mathcal{H}}
\newcommand{\dvc}{d_{\text{VC}}}
\newcommand{\vx}{\mathbf{x}}
\newcommand{\vw}{\mathbf{w}}
\newcommand{\vy}{\mathbf{y}}
\newcommand{\R}{\mathbb{R}}
\newcommand{\T}{^{\top}}
\newcommand{\E}{\mathbb{E}}
\newcommand{\Prob}{\mathbb{P}}
\newcommand{\pts}[1]{\hfill\textbf{(#1~puntos)}}
\newcommand{\linea}[1]{\rule{#1}{0.4pt}}
\newcommand{\ptsinc}[1]{{\unskip\nobreak\hfil\penalty50\hskip1em\hbox{}\nobreak\hfill(#1)\parfillskip=0pt\finalhyphendemerits=0\par}}

\newcommand{\encabezado}[2]{%
\begin{center}
{\large\bfseries Reconocimiento de Patrones}\\[2pt]
Universidad de Sonora \textperiodcentered{} Licenciatura en Ciencias de la Computación\\
Semestre 2027-1 \textperiodcentered{} Prof.\ Julio Waissman Vilanova\\[10pt]
{\Large\bfseries #1}%
\ifdefined\tipodoc\\[4pt]{\large\itshape\tipodoc}\fi
\end{center}
\ifdefined\tipodoc\else
\vspace{6pt}
\textbf{Nombre:} \hrulefill\quad \textbf{Expediente:} \rule{3.5cm}{0.4pt}\par
\vspace{6pt}
\fi
#2
\par\medskip\hrule\medskip
}

\newcommand{\espacio}[1]{\par\vspace{#1}}
\newcommand{\nuevapagina}{\newpage}


\begin{document}

\encabezado{Examen rápido 6 \textperiodcentered{} Redes neuronales}{%
Tiempo: 20 minutos. Responde en esta hoja. Puedes usar calculadora; no se permiten apuntes. Deja indicadas las operaciones.}

\begin{enumerate}

%% ---------------------------------------------------------------
\item \textbf{Hacia atrás en una red diminuta.} Una red con dos entradas, dos neuronas ocultas ReLU y una salida lineal, con pérdida $E = \frac12(\hat y - y)^2$:
\[
\mathbf z^{(1)} = W^{(1)}\vx, \quad \mathbf a^{(1)} = \operatorname{ReLU}\big(\mathbf z^{(1)}\big), \quad \hat y = \mathbf w^{(2)\top}\mathbf a^{(1)} + b^{(2)},
\]
\[
W^{(1)} = \begin{pmatrix} 1 & -1 \\ 0.5 & 0.5 \end{pmatrix}, \qquad \mathbf w^{(2)} = \begin{pmatrix} 2 \\ -1 \end{pmatrix}, \qquad b^{(2)} = 0.5 .
\]
El ejemplo es $\vx = (1, 2)$ con $y = 1$.
\begin{enumerate}
  \item Haz la pasada hacia adelante: $\mathbf z^{(1)}$, $\mathbf a^{(1)}$, $\hat y$ y $E$.
  \espacio{2.5cm}

  \item Haz la pasada hacia atrás: $\delta^{(2)} = \partial E/\partial \hat y$, $\partial E/\partial \mathbf w^{(2)}$, $\partial E / \partial b^{(2)}$, $\boldsymbol\delta^{(1)} = \big(\mathbf w^{(2)}\delta^{(2)}\big) \odot \operatorname{ReLU}'\big(\mathbf z^{(1)}\big)$ y $\partial E/\partial W^{(1)} = \boldsymbol\delta^{(1)}\vx\T$.
  \espacio{4.5cm}

  \item La primera fila de $\partial E/\partial W^{(1)}$ salió en ceros. ¿Por qué? Si esa neurona da $z \le 0$ para \emph{todos} los datos, ¿qué le pasa durante el entrenamiento?
  \espacio{2.2cm}
\end{enumerate}

\nuevapagina
%% ---------------------------------------------------------------
\item \textbf{Contar en una red convolucional.} Imágenes en color de $3 \times 32 \times 32$ (canales $\times$ alto $\times$ ancho) entran a esta red:
\begin{center}
\begin{tabular}{ll}
\toprule
capa & detalle \\ \midrule
conv 1 & 8 filtros de $5 \times 5$, sin relleno, paso 1, ReLU \\
\textit{max-pooling} & $2 \times 2$, paso 2 \\
conv 2 & 16 filtros de $3 \times 3$, relleno 1, paso 2, ReLU \\
lineal & aplanar y 10 salidas \\
\bottomrule
\end{tabular}
\end{center}
Recuerda: $n_{\text{sal}} = \lfloor (n + 2p - k)/s \rfloor + 1$ y $\#\text{parám.} = (C_{\text{ent}}\,k^2 + 1)\,C_{\text{sal}}$.
\begin{enumerate}
  \item Escribe el tamaño de la salida y el número de parámetros de cada capa, y el total.
  \espacio{4cm}

  \item Una capa \emph{densa} que produjera las mismas $8 \times 28 \times 28$ salidas de conv 1 a partir de los $3 \times 32 \times 32$ pixeles, ¿cuántos parámetros tendría? ¿Qué dos ideas explican la diferencia?
  \espacio{2.5cm}

  \item ¿De qué tamaño es el campo receptivo de una neurona de conv 2 (el cuadro de la imagen original que influye en ella)? Usa $r_\ell = r_{\ell-1} + (k-1)\,j_{\ell-1}$, $j_\ell = j_{\ell-1}\,s$, con $r_0 = j_0 = 1$.
  \espacio{2.5cm}
\end{enumerate}

\nuevapagina
%% ---------------------------------------------------------------
\item \textbf{Atención con vectores pequeños.} Una consulta $\mathbf q = (2, 0, 2, 0)$ se compara con tres claves y sus valores ($d_k = 4$):
\[
\mathbf k_1 = (1, 0, 1, 0),\ \ \mathbf k_2 = (0, 1, 0, 1),\ \ \mathbf k_3 = (1, 1, -1, -1); \qquad \mathbf v_1 = (1, 0),\ \ \mathbf v_2 = (0, 1),\ \ \mathbf v_3 = (1, 1).
\]
\begin{enumerate}
  \item Calcula los pesos de atención $\operatorname{softmax}\big(\mathbf q\T\mathbf k_i / \sqrt{d_k}\big)$ y la salida $\sum_i a_i \mathbf v_i$.
  \espacio{3.5cm}

  \item Repite el cálculo de los pesos \emph{sin} dividir entre $\sqrt{d_k}$. ¿Qué cambia? ¿Por qué eso es un problema cuando $d_k$ es grande (por ejemplo, 64)?
  \espacio{2.5cm}

  \item Si la consulta es de la posición 2 en una autoatención con \emph{máscara causal}, y las claves son de las posiciones 1, 2 y 3, ¿cuáles son los pesos y la salida?
  \espacio{2.2cm}
\end{enumerate}

\end{enumerate}


\end{document}
